% =====================================================================
% QA1 - Theoretical Q&A : VR / AR  (bilingual, paracol two-column)
% =====================================================================

\begin{paracol}{2}

\section*{QA1. Theoretical Q\&A --- VR / AR}

\textbf{Q1. What is Virtual Reality (VR)?}\\
\textbf{Virtual Reality} is a computer-generated, fully \textbf{immersive} 3D environment that completely replaces the user's real-world view. The user wears a \textbf{Head-Mounted Display (HMD)} such as the \textbf{Meta Quest 3}, \textbf{HTC Vive Pro}, \textbf{Valve Index}, or \textbf{PlayStation VR2}, plus tracked controllers and sometimes haptic gloves. The system tracks the user's head and hand pose in real time and renders a \textbf{stereoscopic} pair of images (one per eye) at \textbf{90--120 Hz} so the brain perceives true depth. The goals of VR are \textbf{presence} (the feeling of ``being there'') and \textbf{immersion} (sensory replacement). Typical applications include training simulators, surgical rehearsal, exposure therapy, architectural walkthroughs, and games such as \textit{Half-Life: Alyx} and \textit{Beat Saber}.

\switchcolumn

\section*{QA1. Câu hỏi lý thuyết --- VR / AR}

\textbf{Câu 1. Thực tại ảo (VR) là gì?}\\
\textbf{Virtual Reality (VR)} là một môi trường 3D do máy tính tạo ra, \textbf{hoàn toàn thay thế} thế giới thực mà người dùng nhìn thấy. Người dùng đeo một \textbf{Head-Mounted Display (HMD)} như \textbf{Meta Quest 3}, \textbf{HTC Vive Pro}, \textbf{Valve Index} hoặc \textbf{PSVR2}, kèm theo controller và đôi khi haptic gloves. Hệ thống \textbf{tracking} đầu và tay người dùng theo thời gian thực, rồi \textbf{render} một cặp ảnh \textbf{stereoscopic} (mỗi mắt một ảnh) ở \textbf{90--120 Hz} để não cảm nhận chiều sâu. Mục tiêu của VR là tạo \textbf{presence} (cảm giác ``đang ở trong đó'') và \textbf{immersion} (thay thế cảm giác). Ứng dụng phổ biến: huấn luyện, mô phỏng phẫu thuật, trị liệu phơi nhiễm, đi dạo kiến trúc, game như \textit{Half-Life: Alyx}, \textit{Beat Saber}.

\switchcolumn*

\textbf{Q2. What is Augmented Reality (AR)?}\\
\textbf{Augmented Reality} \textit{augments} the real world by overlaying digital content (3D models, text, animations) onto the user's view of the physical environment, instead of replacing it. AR can run on \textbf{smartphones/tablets} (video see-through, e.g.\ Pokemon Go, IKEA Place) or on \textbf{AR glasses} (optical see-through, e.g.\ Microsoft HoloLens 2, Magic Leap 2, Apple Vision Pro in passthrough). Core capabilities are \textbf{6DoF tracking}, \textbf{plane detection}, \textbf{anchoring}, \textbf{occlusion}, and \textbf{light estimation}. Common SDKs are \textbf{ARKit} (iOS), \textbf{ARCore} (Android), \textbf{Vuforia} (cross-platform marker-based), and \textbf{Unity AR Foundation} (unifies them). Applications include navigation arrows on streets, surgical guidance overlay, AR try-on for retail, and industrial maintenance instructions.

\switchcolumn

\textbf{Câu 2. Thực tại tăng cường (AR) là gì?}\\
\textbf{Augmented Reality (AR)} \textit{bổ sung} nội dung số (mô hình 3D, text, animation) lên thế giới thực thay vì thay thế nó. AR chạy trên \textbf{smartphone/tablet} (video see-through, ví dụ Pokemon Go, IKEA Place) hoặc trên \textbf{AR glasses} (optical see-through như HoloLens 2, Magic Leap 2, Apple Vision Pro). Năng lực cốt lõi gồm: \textbf{6DoF tracking}, \textbf{plane detection}, \textbf{anchoring}, \textbf{occlusion}, \textbf{light estimation}. Các SDK phổ biến: \textbf{ARKit} (iOS), \textbf{ARCore} (Android), \textbf{Vuforia} (cross-platform, marker-based), \textbf{Unity AR Foundation} (gói chung). Ứng dụng: chỉ đường ngoài phố, hướng dẫn phẫu thuật, thử đồ ảo, hướng dẫn bảo trì công nghiệp.

\switchcolumn*

\textbf{Q3. What is Mixed Reality (MR)?}\\
\textbf{Mixed Reality} sits between AR and VR on the \textbf{Milgram Reality--Virtuality Continuum}. In MR, virtual objects are not just overlaid but actually \textbf{interact with} the real environment: they understand the geometry of walls and furniture, support \textbf{occlusion} (a real hand can hide a virtual ball), respond to real lighting, and accept physical interaction. MR therefore requires \textbf{spatial mapping} (a real-time mesh of the room produced by depth sensors) plus \textbf{hand and eye tracking}. Devices designed for MR include the \textbf{Microsoft HoloLens 2}, \textbf{Magic Leap 2}, \textbf{Meta Quest 3} (color passthrough), and \textbf{Apple Vision Pro}. Use cases: collaborative engineering review, holographic medical training, factory digital twins.

\switchcolumn

\textbf{Câu 3. Mixed Reality (MR) là gì?}\\
\textbf{Mixed Reality (MR)} nằm giữa AR và VR trên \textbf{Milgram Reality--Virtuality Continuum}. Vật thể ảo không chỉ ``đè lên'' mà thực sự \textbf{tương tác} với môi trường thật: hiểu được hình học tường, bàn ghế, hỗ trợ \textbf{occlusion} (tay thật che được quả bóng ảo), phản ứng với ánh sáng thật, và cho phép tương tác vật lý. MR vì vậy cần \textbf{spatial mapping} (lưới 3D của phòng từ depth sensor) cộng với \textbf{hand tracking} và \textbf{eye tracking}. Thiết bị MR: \textbf{HoloLens 2}, \textbf{Magic Leap 2}, \textbf{Meta Quest 3} (color passthrough), \textbf{Apple Vision Pro}. Use case: review kỹ thuật phối hợp, huấn luyện y khoa hologram, digital twin nhà máy.

\switchcolumn*

\textbf{Q4. What is XR (Extended Reality)?}\\
\textbf{XR} is the \textbf{umbrella term} that encompasses VR, AR, and MR --- ``X'' stands for any reality. It is used in industry standards (e.g., \textbf{OpenXR}) and in research because the boundaries between the three are increasingly blurry: a Quest~3 in passthrough mode is AR/MR, the same device with passthrough off is VR. Talking about ``XR'' lets us describe \textbf{cross-cutting concerns} such as tracking, rendering, input, ergonomics, and content authoring without committing to one mode. The \textbf{OpenXR} standard (Khronos Group) is the de-facto cross-vendor API, supported by Meta, Microsoft, Valve, HTC, Pico, and Apple. Modern game engines like \textbf{Unity} and \textbf{Unreal} target OpenXR so a single codebase runs on all major XR headsets.

\switchcolumn

\textbf{Câu 4. XR (Extended Reality) là gì?}\\
\textbf{XR} là \textbf{thuật ngữ ô} bao gồm VR, AR và MR --- ``X'' nghĩa là bất kỳ loại reality nào. Nó được dùng trong chuẩn công nghiệp (ví dụ \textbf{OpenXR}) và trong nghiên cứu, vì ranh giới giữa ba loại ngày càng mờ: Quest 3 bật passthrough là AR/MR, tắt passthrough thì là VR. Dùng từ ``XR'' giúp ta nói về các vấn đề \textbf{xuyên suốt} như tracking, render, input, ergonomic mà không cần chọn loại cụ thể. Chuẩn \textbf{OpenXR} (Khronos Group) là API cross-vendor de-facto, hỗ trợ bởi Meta, Microsoft, Valve, HTC, Pico, Apple. Các engine \textbf{Unity}, \textbf{Unreal} đều target OpenXR nên một codebase chạy được trên đa số headset.

\switchcolumn*

\textbf{Q5. Compare VR vs AR vs MR (in one paragraph and a table).}\\
The three differ mainly in \textbf{how much real world remains visible} and \textbf{how virtual objects interact with reality}.\\
\begin{tabular}{@{}lccc@{}}
\toprule
 & VR & AR & MR\\
\midrule
Real world & hidden & visible & visible\\
Virtual obj.\ & full env. & overlay & integrated\\
Interaction & with virt. & passive & with real\\
Display & opaque HMD & phone/glass & MR HMD\\
Example & Quest VR & Pokemon Go & HoloLens 2\\
\bottomrule
\end{tabular}\\
In short: \textbf{VR replaces}, \textbf{AR adds on top}, \textbf{MR blends with awareness}. All three sit on the Milgram continuum and are unified under \textbf{XR}.

\switchcolumn

\textbf{Câu 5. So sánh VR, AR, MR (1 đoạn + bảng).}\\
Khác biệt chính nằm ở \textbf{lượng thế giới thực còn nhìn thấy} và \textbf{mức độ tương tác} của vật thể ảo với thực tại.\\
\begin{tabular}{@{}lccc@{}}
\toprule
 & VR & AR & MR\\
\midrule
TG thực & ẩn & thấy & thấy\\
Vật ảo & toàn bộ & overlay & tích hợp\\
Tương tác & với ảo & thụ động & với thực\\
Display & HMD đen & phone/kính & MR HMD\\
Ví dụ & Quest VR & Pokemon Go & HoloLens 2\\
\bottomrule
\end{tabular}\\
Tóm lại: \textbf{VR thay thế}, \textbf{AR thêm vào}, \textbf{MR hòa quyện}. Cả ba nằm trên Milgram continuum và gộp chung thành \textbf{XR}.

\switchcolumn*

\textbf{Q6. What is an HMD and what are its main components?}\\
A \textbf{Head-Mounted Display (HMD)} is the wearable display that puts visuals close to the eyes. Its main components are: (1) two \textbf{micro-displays} (LCD, OLED, or micro-OLED), one per eye, with resolutions like $2064\times2208$ per eye on Quest 3; (2) \textbf{Fresnel} or \textbf{pancake lenses} that magnify the image so the eye can focus at near distance; (3) an \textbf{IPD adjustment} for inter-pupillary distance; (4) an \textbf{IMU} (gyroscope + accelerometer + magnetometer) for fast head rotation tracking; (5) outward-facing \textbf{cameras} for inside-out 6DoF tracking and passthrough; (6) sometimes \textbf{eye-tracking cameras} for foveated rendering. Standalone HMDs (Quest 3) include an SoC like Snapdragon XR2~Gen~2; PC-VR HMDs (Vive Pro 2, Index) stream from a GPU through DisplayPort.

\switchcolumn

\textbf{Câu 6. HMD là gì? Các thành phần chính.}\\
\textbf{Head-Mounted Display (HMD)} là thiết bị đeo đầu đưa hình ảnh sát mắt. Thành phần chính: (1) hai \textbf{micro-display} (LCD, OLED hoặc micro-OLED), mỗi mắt một, ví dụ Quest 3 có $2064\times2208$ mỗi mắt; (2) \textbf{Fresnel} hoặc \textbf{pancake lens} phóng đại để mắt focus được ở khoảng gần; (3) \textbf{IPD adjustment} để chỉnh khoảng cách hai mắt; (4) \textbf{IMU} (gyro + accel + magnetometer) đo xoay đầu siêu nhanh; (5) \textbf{camera} hướng ngoài để inside-out 6DoF tracking và passthrough; (6) đôi khi có \textbf{eye-tracking camera} cho foveated rendering. HMD standalone (Quest 3) có SoC kiểu Snapdragon XR2 Gen 2; PC-VR HMD (Vive Pro 2, Index) nhận stream từ GPU qua DisplayPort.

\switchcolumn*

\textbf{Q7. What is 6DoF tracking? Compare with 3DoF.}\\
\textbf{Degrees of Freedom (DoF)} = number of independent motions the system can sense. \textbf{3DoF} senses only \textit{rotations} around three axes (yaw, pitch, roll) --- the user can look around but cannot walk; typical of cheap mobile VR like the original Cardboard or Oculus Go. \textbf{6DoF} adds the three \textit{translations} (forward/back, left/right, up/down), so the user can physically walk around and lean. 6DoF requires either \textbf{outside-in} tracking (e.g., HTC Vive Lighthouse base stations emitting IR sweeps) or \textbf{inside-out} tracking (cameras on the HMD doing visual SLAM, e.g., Quest 3, HoloLens 2). 6DoF is essential for \textbf{presence}, room-scale interaction, and is the standard for any modern VR/AR experience.

\switchcolumn

\textbf{Câu 7. 6DoF tracking là gì? So với 3DoF?}\\
\textbf{Degrees of Freedom (DoF)} = số chuyển động độc lập hệ thống cảm nhận được. \textbf{3DoF} chỉ cảm nhận \textit{xoay} quanh 3 trục (yaw/pitch/roll) --- người dùng nhìn quanh được nhưng không đi được; đặc trưng của mobile VR rẻ như Cardboard hoặc Oculus Go. \textbf{6DoF} thêm 3 chiều \textit{tịnh tiến} (trước/sau, trái/phải, lên/xuống), người dùng đi và nghiêng được. 6DoF cần \textbf{outside-in} (HTC Vive Lighthouse phát IR sweep) hoặc \textbf{inside-out} (camera trên HMD chạy visual SLAM, ví dụ Quest 3, HoloLens 2). 6DoF là điều kiện cần để có \textbf{presence}, room-scale interaction, và là chuẩn của VR/AR hiện đại.

\switchcolumn*

\textbf{Q8. What is SLAM and how does it work?}\\
\textbf{SLAM = Simultaneous Localization And Mapping}. The headset (or robot, drone, phone) builds a map of an unknown environment \textit{while} tracking its own pose inside that map, in real time. Pipeline: (1) \textbf{Feature extraction} from each frame (corners/keypoints via FAST, ORB, or learned features); (2) \textbf{Feature matching} between consecutive frames; (3) \textbf{Pose estimation} via PnP / essential matrix; (4) \textbf{IMU fusion} with an Extended Kalman Filter or factor graph for smoothness during fast motion; (5) \textbf{Mapping}: triangulate matched features into 3D landmarks; (6) \textbf{Loop closure} corrects accumulated \textbf{drift} when a previously seen place is recognized; (7) \textbf{Bundle adjustment} optimizes all poses + landmarks. Visual--inertial SLAM (VI-SLAM) underlies ARKit, ARCore, HoloLens, and Quest tracking.

\switchcolumn

\textbf{Câu 8. SLAM là gì? Hoạt động ra sao?}\\
\textbf{SLAM = Simultaneous Localization And Mapping}. Headset (hoặc robot, drone, phone) vừa \textit{xây bản đồ} môi trường lạ, vừa \textit{định vị bản thân} trong bản đồ đó, theo thời gian thực. Pipeline: (1) \textbf{Feature extraction} từ mỗi frame (góc/keypoint qua FAST, ORB hoặc deep features); (2) \textbf{Feature matching} giữa các frame liên tiếp; (3) \textbf{Pose estimation} qua PnP / essential matrix; (4) \textbf{IMU fusion} bằng Extended Kalman Filter hoặc factor graph cho mượt khi chuyển động nhanh; (5) \textbf{Mapping}: tam giác hóa các feature thành landmark 3D; (6) \textbf{Loop closure} sửa \textbf{drift} khi nhận ra nơi cũ; (7) \textbf{Bundle adjustment} tối ưu toàn bộ pose + landmark. Visual--inertial SLAM (VI-SLAM) là nền của ARKit, ARCore, HoloLens, Quest.

\switchcolumn*

\textbf{Q9. What is presence? What is immersion?}\\
\textbf{Immersion} is an \textbf{objective property of the system}: how well the hardware replaces sensory input (high resolution, wide \textbf{FOV}, low latency, spatial audio, haptics). \textbf{Presence} is the \textbf{subjective feeling} of ``being there'' that the user experiences \textit{because} immersion is high. Immersion enables presence but does not guarantee it: bad content, low frame rate, or cybersickness can break presence even on great hardware. Presence is usually measured with \textbf{questionnaires} like the Slater--Usoh--Steed (SUS) or Witmer--Singer Presence Questionnaire, or with physiological proxies (heart rate, EDA). Designers maximize presence with consistent \textbf{6DoF tracking}, $\geq$~90~Hz refresh, $<$~20~ms motion-to-photon latency, spatial audio, and believable interaction.

\switchcolumn

\textbf{Câu 9. Presence là gì? Immersion là gì?}\\
\textbf{Immersion} là \textbf{thuộc tính khách quan của hệ thống}: phần cứng thay thế giác quan tốt đến đâu (độ phân giải cao, \textbf{FOV} rộng, latency thấp, âm thanh không gian, haptic). \textbf{Presence} là \textbf{cảm giác chủ quan} ``đang ở đó'' mà người dùng cảm thấy \textit{nhờ} immersion cao. Immersion là điều kiện cần nhưng chưa đủ: nội dung dở, FPS thấp, cybersickness đều có thể phá vỡ presence. Presence thường đo bằng \textbf{questionnaire} như Slater--Usoh--Steed (SUS), Witmer--Singer, hoặc đo sinh lý (nhịp tim, EDA). Thiết kế tốt để tối đa presence: \textbf{6DoF tracking} ổn định, $\geq$~90~Hz, motion-to-photon $<$~20~ms, spatial audio, tương tác thuyết phục.

\switchcolumn*

\textbf{Q10. How does VR work step-by-step (the rendering loop)?}\\
The VR loop runs at 90--120~Hz and must finish each frame in $\sim$11~ms or less:\\
\textbf{(1) Sense:} read \textbf{IMU} + cameras to update head pose; \textbf{(2) Predict:} extrapolate the pose to the photon-display time (\textbf{motion prediction}); \textbf{(3) Update:} run game logic, physics, animation; \textbf{(4) Cull:} frustum + occlusion culling; \textbf{(5) Render twice:} produce the \textit{left-eye} and \textit{right-eye} images using two slightly offset projection matrices (one per camera, separated by IPD); \textbf{(6) Distortion correction:} pre-warp each image to cancel the lens's pincushion distortion and chromatic aberration; \textbf{(7) Asynchronous Time Warp / Space Warp:} just before scan-out, re-project the finished image with the latest pose to compensate for any remaining latency; \textbf{(8) Scan out} to the panel, scanline-synced. Spatial audio is computed in parallel with HRTFs.

\switchcolumn

\textbf{Câu 10. VR hoạt động từng bước thế nào (rendering loop)?}\\
Vòng VR chạy 90--120~Hz, mỗi frame phải xong trong $\sim$11~ms:\\
\textbf{(1) Sense:} đọc \textbf{IMU} + camera, cập nhật head pose; \textbf{(2) Predict:} ngoại suy pose tới thời điểm photon ra màn (\textbf{motion prediction}); \textbf{(3) Update:} game logic, physics, animation; \textbf{(4) Cull:} frustum + occlusion culling; \textbf{(5) Render hai lần:} ảnh \textit{mắt trái} và \textit{mắt phải} với hai ma trận chiếu lệch nhau (cách nhau IPD); \textbf{(6) Distortion correction:} pre-warp để khử pincushion và chromatic aberration của thấu kính; \textbf{(7) Asynchronous Time Warp / Space Warp:} ngay trước scan-out, re-project ảnh đã render bằng pose mới nhất để bù latency; \textbf{(8) Scan out} ra panel theo scanline. Spatial audio tính song song bằng HRTF.

\switchcolumn*

\textbf{Q11. How does the AR pipeline work step-by-step?}\\
\textbf{(1) Camera capture:} grab the next frame from the rear camera (or passthrough cameras); \textbf{(2) Tracking:} run \textbf{VI-SLAM} to estimate the device's 6DoF pose; \textbf{(3) Scene understanding:} detect \textbf{planes} (floor, table), feature points, meshes, faces, hands, image targets; \textbf{(4) Light estimation:} infer ambient intensity, color, and a directional source (or spherical harmonics) from the camera image; \textbf{(5) Anchor management:} attach virtual objects to anchors that are stable in the real world; \textbf{(6) Rendering:} render virtual objects with the camera's intrinsics and the estimated pose; apply real-world lighting; perform \textbf{occlusion} using the depth map or scene mesh; \textbf{(7) Composite:} blend the rendered overlay with the camera image (video see-through) or display on the optical combiner (optical see-through); \textbf{(8) Display} at $\geq$~60~Hz.

\switchcolumn

\textbf{Câu 11. AR pipeline hoạt động từng bước?}\\
\textbf{(1) Camera capture:} lấy frame mới từ camera sau (hoặc camera passthrough); \textbf{(2) Tracking:} chạy \textbf{VI-SLAM} để ước lượng 6DoF pose của thiết bị; \textbf{(3) Scene understanding:} dò \textbf{plane} (sàn, bàn), feature point, mesh, mặt người, tay, image target; \textbf{(4) Light estimation:} suy ra cường độ, màu sắc môi trường + nguồn sáng định hướng (hoặc spherical harmonics) từ ảnh camera; \textbf{(5) Anchor management:} gắn vật ảo vào anchor cố định trong không gian thật; \textbf{(6) Rendering:} render vật ảo theo intrinsic của camera và pose ước lượng; áp ánh sáng thực; thực hiện \textbf{occlusion} dùng depth map hoặc scene mesh; \textbf{(7) Composite:} blend overlay với ảnh camera (video see-through) hoặc hiển thị qua combiner quang (optical see-through); \textbf{(8) Hiển thị} ở $\geq$~60~Hz.

\switchcolumn*

\textbf{Q12. How does inside-out tracking work?}\\
In \textbf{inside-out tracking}, the cameras are \textit{on the headset} and look outward; there are no external base stations. Each frame, the system extracts \textbf{visual features} (corners), matches them across multiple cameras (stereo or wider arrays) and across time, and combines this with the \textbf{IMU} via \textbf{visual--inertial odometry}. The result is a 6DoF pose at very high rate (1000~Hz from IMU, corrected by 30--60~Hz visual updates). For controllers, the same outward cameras observe \textbf{IR LEDs} on the controller (Quest's constellation) or perform fully optical hand tracking (HoloLens, Quest hand tracking, Vision Pro). Pros: no setup, portable, fewer cables. Cons: needs textured environment, can lose tracking in the dark or in front of a blank wall, controller can be lost behind the back.

\switchcolumn

\textbf{Câu 12. Inside-out tracking hoạt động ra sao?}\\
Trong \textbf{inside-out tracking}, camera nằm \textit{trên headset}, hướng ra ngoài; không cần base station. Mỗi frame, hệ thống trích \textbf{feature} (góc), match qua nhiều camera (stereo hoặc nhiều hơn) và qua thời gian, kết hợp với \textbf{IMU} bằng \textbf{visual--inertial odometry}. Kết quả là 6DoF pose tần số rất cao (IMU 1000 Hz, hiệu chỉnh bằng visual 30--60 Hz). Với controller: cùng camera đó quan sát \textbf{IR LED} trên tay cầm (Quest constellation) hoặc làm hand tracking thuần optical (HoloLens, Quest hand tracking, Vision Pro). Ưu: không cần setup, di động, ít cáp. Nhược: cần môi trường có texture, mất tracking trong tối hoặc trước tường trắng, controller có thể mất khi giấu sau lưng.

\switchcolumn*

\textbf{Q13. Compare outside-in (Lighthouse) vs inside-out tracking.}\\
\textbf{Outside-in} uses external sensors (Lighthouse base stations sweeping IR planes; or external cameras like the original Oculus Rift CV1) that observe the headset/controllers; \textbf{inside-out} reverses this: cameras are on the headset. Outside-in offers \textbf{sub-millimeter precision}, near-zero jitter, large play areas (Vive Lighthouse covers $\sim$10~m), and reliable controller tracking even behind the back. Its drawbacks: cumbersome setup, fixed play space, needs power outlets and line-of-sight. Inside-out is \textbf{plug-and-play}, portable, cheaper, supports any room, but has slightly more jitter, can drift, struggles in low light or feature-poor rooms, and may lose controllers occluded from the headset cameras. Modern consumer HMDs (Quest, HoloLens, Vision Pro, PSVR2) all use inside-out; pro/research setups still favor Lighthouse.

\switchcolumn

\textbf{Câu 13. So sánh outside-in (Lighthouse) vs inside-out.}\\
\textbf{Outside-in} dùng cảm biến ngoài (Lighthouse phát mặt phẳng IR; hoặc camera ngoài như Rift CV1) quan sát headset/controller; \textbf{inside-out} ngược lại, camera nằm trên headset. Outside-in cho \textbf{độ chính xác sub-mm}, gần như không jitter, vùng chơi rộng (Vive Lighthouse phủ $\sim$10~m), tracking controller ổn định kể cả sau lưng. Nhược: setup phức tạp, vùng chơi cố định, cần ổ cắm và line-of-sight. Inside-out là \textbf{plug-and-play}, di động, rẻ, hỗ trợ mọi phòng, nhưng jitter hơn, có drift, kém trong tối hoặc phòng ít texture, mất controller khi bị che. HMD tiêu dùng hiện nay (Quest, HoloLens, Vision Pro, PSVR2) dùng inside-out; setup chuyên nghiệp/nghiên cứu vẫn ưa Lighthouse.

\switchcolumn*

\textbf{Q14. How does marker-based AR work?}\\
\textbf{Marker-based AR} relies on a known image (\textbf{fiducial marker} like ARTag/ArUco, or a textured \textbf{image target} in Vuforia) printed in the real world. Steps: (1) The camera grabs a frame; (2) The image is binarized/thresholded and contours are found; (3) Quadrilateral contours are matched against the marker's known pattern (or feature descriptors are matched against the image database); (4) The four corners give a 2D--3D correspondence; (5) Solving \textbf{PnP (Perspective-n-Point)} with the camera intrinsics yields the 6DoF pose of the marker relative to the camera; (6) A virtual object is rendered at that pose. Marker-based AR is \textbf{robust, fast, and free of drift}, used in print-to-AR books, museum exhibits, packaging promotions, and education --- but the marker must remain visible.

\switchcolumn

\textbf{Câu 14. Marker-based AR hoạt động thế nào?}\\
\textbf{Marker-based AR} dựa vào ảnh đã biết (\textbf{fiducial marker} kiểu ARTag/ArUco, hoặc \textbf{image target} có texture trong Vuforia) đặt trong thế giới thực. Bước: (1) Camera lấy frame; (2) Ảnh được binarize/threshold, tìm contour; (3) Contour tứ giác được match với pattern marker (hoặc feature descriptor match với database); (4) Bốn góc cho tương ứng 2D--3D; (5) Giải \textbf{PnP (Perspective-n-Point)} với intrinsic camera $\Rightarrow$ 6DoF pose của marker so với camera; (6) Render vật ảo tại pose đó. Marker-based AR \textbf{ổn định, nhanh, không drift}, dùng trong sách AR, bảo tàng, bao bì khuyến mại, giáo dục --- nhưng marker phải luôn nhìn thấy.

\switchcolumn*

\textbf{Q15. Compare marker-based vs markerless AR.}\\
\textbf{Marker-based} requires a printed marker/image target in the scene; \textbf{markerless} relies on natural features and SLAM. Marker-based is \textbf{simple, robust, drift-free, low-cost} but constrained: the experience dies if the marker is occluded, lit poorly, or out of frame. Markerless (ARKit, ARCore, HoloLens) discovers planes and feature points anywhere, supports walking around the scene, and enables persistent \textbf{world anchors} via cloud anchoring (ARCore Cloud Anchors, ARKit ARWorldMap). Its trade-offs are higher compute, dependence on texture and lighting, and \textbf{drift} that needs loop closure. Practical rule: use marker-based for short, controlled experiences (museum cards, packaging, education); use markerless for room-scale, navigation, and large-area AR.

\switchcolumn

\textbf{Câu 15. So sánh marker-based vs markerless AR.}\\
\textbf{Marker-based} cần marker/image target in sẵn; \textbf{markerless} dựa vào feature tự nhiên + SLAM. Marker-based \textbf{đơn giản, ổn, không drift, rẻ} nhưng bị ràng buộc: marker bị che, thiếu sáng, ngoài khung là ``chết''. Markerless (ARKit, ARCore, HoloLens) phát hiện plane và feature mọi nơi, đi lại được, hỗ trợ \textbf{world anchor} bền qua cloud (ARCore Cloud Anchors, ARKit ARWorldMap). Nhược điểm: tốn compute, phụ thuộc texture và ánh sáng, có \textbf{drift} cần loop closure. Quy tắc thực tế: marker-based cho trải nghiệm ngắn, có kiểm soát (bảo tàng, bao bì, giáo dục); markerless cho room-scale, dẫn đường, AR diện rộng.

\switchcolumn*

\textbf{Q16. Compare optical see-through vs video see-through.}\\
\textbf{Optical see-through (OST)}: light from the real world reaches the eye through a transparent \textbf{combiner} (waveguide, beam-splitter); virtual content is projected onto it. Used by HoloLens 2, Magic Leap 2, smart glasses. Pros: \textbf{zero camera latency} for the real view, natural depth, safer (you actually see). Cons: limited \textbf{FOV} ($\sim$50$^\circ$ on HL2), virtual objects look semi-transparent, hard to render true black, expensive optics. \textbf{Video see-through (VST)}: outward cameras capture the world, the system composites virtuals on top, and the result is shown on opaque displays (Quest 3, Vision Pro, Varjo XR-4). Pros: \textbf{wide FOV}, perfect occlusion, dynamic range, full opacity for virtuals. Cons: camera latency, distortion, lower fidelity than direct vision, total visual blackout if the system crashes.

\switchcolumn

\textbf{Câu 16. So sánh optical vs video see-through.}\\
\textbf{Optical see-through (OST)}: ánh sáng thực đi qua \textbf{combiner} trong suốt (waveguide, beam-splitter); nội dung ảo được chiếu lên đó. HoloLens 2, Magic Leap 2, smart glasses dùng cách này. Ưu: \textbf{0 ms camera latency} cho cảnh thực, depth tự nhiên, an toàn hơn. Nhược: \textbf{FOV} hẹp ($\sim$50$^\circ$ ở HL2), vật ảo trong mờ, khó render màu đen thật, optics đắt. \textbf{Video see-through (VST)}: camera ngoài quay cảnh thực, ghép ảo lên, hiển thị qua màn hình kín (Quest 3, Vision Pro, Varjo XR-4). Ưu: \textbf{FOV rộng}, occlusion hoàn hảo, dải động cao, ảo phủ kín được. Nhược: latency camera, distortion, độ trung thực kém mắt thường, mất hình toàn bộ nếu hệ thống crash.

\switchcolumn*

\textbf{Q17. How does stereoscopic rendering produce depth?}\\
The brain extracts depth from many cues; the strongest near-field cue is \textbf{binocular disparity} --- the slight horizontal offset between the images received by the left and right eye. To replicate this, a VR engine maintains \textbf{two virtual cameras} separated by the user's \textbf{IPD} ($\sim$63~mm) and renders each frame twice with slightly different view matrices but the same model and projection setup (with appropriate eye offsets). Each rendered image is sent to its eye's panel through the lens. The brain fuses the two images and infers depth. Additional cues are \textbf{motion parallax} (provided by 6DoF head tracking) and \textbf{vergence}. Mismatch between vergence and the fixed lens focus causes the \textbf{vergence--accommodation conflict (VAC)}, a known cause of fatigue and a target of varifocal R\&D.

\switchcolumn

\textbf{Câu 17. Stereoscopic render tạo depth thế nào?}\\
Não suy ra depth từ nhiều cue; cue gần mạnh nhất là \textbf{binocular disparity} --- chênh lệch ngang giữa ảnh hai mắt. Để mô phỏng, engine VR đặt \textbf{hai camera ảo} cách nhau bằng \textbf{IPD} ($\sim$63~mm), mỗi frame render hai lần với view matrix lệch (chung model + projection có offset mắt). Mỗi ảnh đi qua thấu kính tới panel của mắt tương ứng. Não fuse hai ảnh và tính ra depth. Cue phụ: \textbf{motion parallax} (do 6DoF head tracking cung cấp) và \textbf{vergence}. Khi vergence lệch tiêu cự cố định của thấu kính sẽ gây \textbf{vergence--accommodation conflict (VAC)} --- nguyên nhân mỏi mắt, đang là chủ đề R\&D varifocal.

\switchcolumn*

\textbf{Q18. How does foveated rendering work and why use it?}\\
The human retina has high resolution only in the \textbf{fovea} ($\sim$5$^\circ$ at the center of gaze); peripheral vision is much lower resolution. \textbf{Foveated rendering} exploits this by rendering the \textbf{gaze-centered region at full resolution} and the periphery at progressively lower resolution, then blending. \textbf{Eye-tracked} variant (PSVR2, Vision Pro, Quest Pro) uses real-time gaze; \textbf{fixed} variant (Quest 2/3 default) just lowers resolution at panel edges. Benefits: \textbf{30--70\% GPU savings}, enabling higher-fidelity scenes or lower power on standalone HMDs. Implementation in graphics APIs: \textbf{Variable Rate Shading (VRS)}, Vulkan VK\_EXT\_fragment\_density\_map, NVIDIA VRSS. The technique is essential for high-resolution headsets ($>$~3K per eye) where full-res rendering would not fit a 90~Hz budget.

\switchcolumn

\textbf{Câu 18. Foveated rendering hoạt động thế nào, vì sao dùng?}\\
Võng mạc người chỉ có độ phân giải cao ở \textbf{fovea} ($\sim$5$^\circ$ trung tâm); vùng ngoại vi thấp hơn nhiều. \textbf{Foveated rendering} tận dụng điều đó: render \textbf{vùng nhìn} ở full resolution, vùng ngoại vi giảm dần resolution rồi blend. Bản \textbf{eye-tracked} (PSVR2, Vision Pro, Quest Pro) dùng gaze thực; bản \textbf{fixed} (Quest 2/3 mặc định) chỉ giảm resolution ở rìa panel. Lợi ích: \textbf{tiết kiệm GPU 30--70\%}, cho phép scene đẹp hơn hoặc tiết kiệm pin trên HMD standalone. Triển khai bằng \textbf{Variable Rate Shading (VRS)}, Vulkan VK\_EXT\_fragment\_density\_map, NVIDIA VRSS. Kỹ thuật này gần như bắt buộc với headset $>$~3K mỗi mắt vì render full-res không kịp 90~Hz.

\switchcolumn*

\textbf{Q19. Why is low latency critical in VR? What is motion-to-photon?}\\
\textbf{Motion-to-photon (M2P) latency} is the total time from a real head movement until the corresponding photon hits the eye. If M2P is too high, the image lags the head, the brain detects a mismatch between vestibular and visual systems, and \textbf{cybersickness} results. Research and industry consensus: \textbf{M2P should be under 20~ms}; great systems achieve \textbf{$<$~15~ms}. Latency is reduced by: (1) high-frequency IMU + prediction; (2) pipelined GPU rendering; (3) \textbf{Asynchronous Time Warp (ATW)} re-projects the last rendered frame using fresh head pose right before scan-out; (4) \textbf{Asynchronous Space Warp (ASW)} synthesizes intermediate frames; (5) racing-the-beam scan-out. High refresh rate (90/120/144~Hz) further reduces effective latency and judder.

\switchcolumn

\textbf{Câu 19. Vì sao latency thấp quan trọng trong VR? Motion-to-photon là gì?}\\
\textbf{Motion-to-photon (M2P) latency} là tổng thời gian từ khi đầu chuyển động thực tới khi photon tới mắt. Nếu M2P cao, ảnh trễ so với đầu, não phát hiện lệch giữa tiền đình và thị giác $\Rightarrow$ \textbf{cybersickness}. Đồng thuận: \textbf{M2P dưới 20~ms}; hệ tốt đạt \textbf{$<$~15~ms}. Cách giảm latency: (1) IMU tần số cao + prediction; (2) GPU render pipelined; (3) \textbf{Asynchronous Time Warp (ATW)} re-project frame cũ bằng pose mới trước khi scan-out; (4) \textbf{Asynchronous Space Warp (ASW)} sinh frame trung gian; (5) racing-the-beam scan-out. Refresh rate cao (90/120/144~Hz) giảm thêm latency hữu hiệu và judder.

\switchcolumn*

\textbf{Q20. Why does motion sickness (cybersickness) occur in VR? How to prevent it?}\\
\textbf{Cybersickness} arises mainly from a \textbf{sensory conflict} between the visual system (which says ``I'm moving'') and the \textbf{vestibular system} (which says ``I'm sitting still''). Triggers: high latency, low frame rate, artificial \textbf{smooth locomotion}, vection, narrow FOV vignetting absent, miscalibrated IPD, lens distortion errors. Mitigation: \textbf{(1)} keep M2P $<$~20~ms and frame rate $\geq$~90~Hz; \textbf{(2)} use \textbf{teleport locomotion} or short \textbf{snap turns} (30$^\circ$); \textbf{(3)} apply \textbf{tunnel-vision FOV reducer} during artificial movement; \textbf{(4)} provide a \textbf{static reference} like a cockpit or nose anchor; \textbf{(5)} keep the floor calibrated and the horizon stable; \textbf{(6)} let users sit and start with short sessions; \textbf{(7)} avoid acceleration --- use constant velocity. Always include comfort settings in shipping VR apps.

\switchcolumn

\textbf{Câu 20. Vì sao bị say (cybersickness) trong VR? Cách phòng?}\\
\textbf{Cybersickness} chủ yếu do \textbf{xung đột giác quan} giữa thị giác (``đang di chuyển'') và \textbf{tiền đình} (``đang ngồi yên''). Nguyên nhân: latency cao, FPS thấp, \textbf{smooth locomotion} nhân tạo, vection, không có vignette, IPD sai, distortion lệch. Cách giảm: \textbf{(1)} giữ M2P $<$~20~ms và $\geq$~90~Hz; \textbf{(2)} dùng \textbf{teleport} hoặc \textbf{snap turn} 30$^\circ$; \textbf{(3)} \textbf{tunnel-vision} thu hẹp FOV khi di chuyển nhân tạo; \textbf{(4)} thêm \textbf{tham chiếu tĩnh} (cockpit, nose anchor); \textbf{(5)} giữ sàn calibrated, đường chân trời ổn định; \textbf{(6)} ngồi và bắt đầu phiên ngắn; \textbf{(7)} tránh acceleration, dùng vận tốc đều. Ứng dụng VR thương mại bắt buộc có comfort settings.

\switchcolumn*

\textbf{Q21. Why is high refresh rate (90+ Hz) needed?}\\
At lower refresh rates (60~Hz or below), persistent flicker and judder are visible because each frame is held in front of the eye while the head moves; this causes \textbf{smearing} and contributes strongly to nausea. \textbf{90~Hz is the empirical comfort threshold}; \textbf{120~Hz} is preferred for fast titles (Beat Saber, racing) and \textbf{144~Hz} on Index. Higher rates also lower effective M2P latency by reducing the wait until the next vsync. To hit these rates, engines use \textbf{stereo instancing / multiview} (rendering both eyes in one pass), \textbf{single-pass instanced rendering}, fixed/eye-tracked foveation, lower shader complexity, and \textbf{ATW/ASW} as a safety net when rendering misses the deadline. Standalone HMDs additionally use dynamic resolution and aggressive LOD to maintain rate.

\switchcolumn

\textbf{Câu 21. Vì sao cần refresh rate cao (90+ Hz)?}\\
Ở refresh thấp (60~Hz trở xuống), flicker và judder rất rõ vì mỗi frame ``đứng yên'' trước mắt trong khi đầu di chuyển, gây \textbf{smearing} và buồn nôn. \textbf{90~Hz là ngưỡng thoải mái thực nghiệm}; \textbf{120~Hz} cho game nhanh (Beat Saber, racing); \textbf{144~Hz} trên Index. Refresh cao cũng giảm M2P do thời gian chờ vsync ngắn. Để đạt được, engine dùng \textbf{stereo instancing/multiview} (render hai mắt trong một pass), \textbf{single-pass instanced rendering}, foveation cố định/eye-tracked, shader nhẹ, \textbf{ATW/ASW} làm lưới an toàn khi miss deadline. HMD standalone còn dùng dynamic resolution và LOD để giữ FPS.

\switchcolumn*

\textbf{Q22. What is FOV and why does it matter? What is IPD?}\\
\textbf{FOV (Field of View)} is the angular size of what the user can see. Human binocular FOV is $\sim$210$^\circ$ horizontally; consumer VR HMDs offer 90$^\circ$--110$^\circ$ (Quest 3 $\sim$110$^\circ$, Index $\sim$130$^\circ$, Varjo XR-4 $\sim$120$^\circ$ at higher resolution; HoloLens 2 only $\sim$52$^\circ$). \textbf{Wider FOV $\Rightarrow$ stronger immersion} but more pixels and lens complexity. \textbf{IPD (Inter-Pupillary Distance)} is the distance between the user's pupils, typically 58--72~mm; HMDs need to match the user's IPD by mechanical or software adjustment for a sharp, comfortable image and correct stereo. Wrong IPD causes blur, eye strain, and incorrect depth scale. Eye-tracking can perform \textbf{automatic IPD compensation} in software.

\switchcolumn

\textbf{Câu 22. FOV là gì, vì sao quan trọng? IPD là gì?}\\
\textbf{FOV (Field of View)} là góc thị trường người dùng thấy. Mắt người: $\sim$210$^\circ$ ngang; VR tiêu dùng: 90--110$^\circ$ (Quest 3 $\sim$110$^\circ$, Index $\sim$130$^\circ$, Varjo XR-4 $\sim$120$^\circ$ độ phân giải cao; HoloLens 2 chỉ $\sim$52$^\circ$). \textbf{FOV rộng $\Rightarrow$ immersion mạnh} nhưng tốn pixel và optics phức tạp. \textbf{IPD (Inter-Pupillary Distance)} là khoảng cách hai đồng tử, 58--72~mm; HMD phải khớp IPD bằng cơ học hoặc phần mềm để ảnh nét và stereo đúng. IPD sai $\Rightarrow$ mờ, mỏi mắt, sai tỉ lệ depth. Eye-tracking có thể \textbf{IPD compensation} tự động bằng phần mềm.

\switchcolumn*

\textbf{Q23. What is the screen-door effect? How is it solved?}\\
The \textbf{screen-door effect (SDE)} is the appearance of dark gaps between subpixels --- the user sees through a fine ``mesh'' over the image. It happens because lenses magnify the panel close to the eye, exaggerating the gaps in the \textbf{PenTile / RGB stripe} subpixel layout. SDE is reduced by: (1) higher \textbf{pixel density} (PPD --- pixels per degree); (2) \textbf{full-RGB stripe} instead of PenTile; (3) \textbf{micro-OLED} panels that pack subpixels tightly (Vision Pro, Bigscreen Beyond); (4) \textbf{pancake lenses} that allow finer optical performance; (5) optical anti-aliasing diffusers. Modern HMDs (Quest 3 with $\sim$25 PPD, Vision Pro with $\sim$34 PPD) have nearly eliminated SDE; older HMDs (Vive, Rift CV1) had it badly.

\switchcolumn

\textbf{Câu 23. Screen-door effect là gì? Giải quyết ra sao?}\\
\textbf{Screen-door effect (SDE)} là hiện tượng nhìn thấy khoảng tối giữa các subpixel --- ảnh như có ``lưới''. Lý do: thấu kính phóng đại panel sát mắt, làm rõ khe của subpixel layout \textbf{PenTile / RGB stripe}. Cách giảm: (1) tăng \textbf{pixel density} (PPD); (2) dùng \textbf{full-RGB stripe} thay PenTile; (3) panel \textbf{micro-OLED} đóng subpixel rất khít (Vision Pro, Bigscreen Beyond); (4) \textbf{pancake lens} cho optics mịn hơn; (5) diffuser quang chống aliasing. HMD mới (Quest 3 $\sim$25 PPD, Vision Pro $\sim$34 PPD) gần như hết SDE; HMD cũ (Vive, Rift CV1) bị rõ.

\switchcolumn*

\textbf{Q24. Compare HMD display technologies: LCD, OLED, micro-OLED, waveguide.}\\
\textbf{LCD}: cheap, high pixel density, full-RGB stripe, but limited contrast, gray blacks, slower response, backlight bleed. Used in Quest 2/3, Pico 4, Varjo. \textbf{OLED}: per-pixel emission, deep blacks, fast response, but typically PenTile (worse SDE), risk of burn-in. Used in Vive Pro, PSVR2 (HDR OLED). \textbf{Micro-OLED} (silicon-backed OLED): tiny panel, very high pixel density ($>$~3000~PPI), excellent blacks, low power; used in Vision Pro (4K per eye), Bigscreen Beyond, Varjo XR-4. \textbf{Waveguides} (used in HoloLens 2, Magic Leap 2): not a panel but a light-guide that channels micro-display output into the eye through diffraction gratings; enables thin AR glasses but limits FOV and adds rainbow artifacts.

\switchcolumn

\textbf{Câu 24. So sánh display HMD: LCD, OLED, micro-OLED, waveguide.}\\
\textbf{LCD}: rẻ, mật độ pixel cao, full-RGB stripe, nhưng tương phản kém, đen ngả xám, chậm, lọt sáng. Dùng ở Quest 2/3, Pico 4, Varjo. \textbf{OLED}: phát sáng từng pixel, đen sâu, đáp ứng nhanh, nhưng thường PenTile (SDE rõ hơn), nguy cơ burn-in. Dùng ở Vive Pro, PSVR2 (HDR OLED). \textbf{Micro-OLED} (OLED trên silicon): panel rất nhỏ, mật độ rất cao ($>$~3000 PPI), đen tốt, ít điện; dùng ở Vision Pro (4K mỗi mắt), Bigscreen Beyond, Varjo XR-4. \textbf{Waveguide} (HoloLens 2, Magic Leap 2): không phải panel mà là ống dẫn sáng đưa output micro-display vào mắt qua cách tử nhiễu xạ; cho kính AR mỏng nhưng FOV hẹp và có rainbow artifact.

\switchcolumn*

\textbf{Q25. What are the main components of a VR/AR system?}\\
A complete XR system has: \textbf{(1) Display}: the HMD (panels + lenses); \textbf{(2) Tracking}: 6DoF head tracking (IMU + cameras), controller tracking, hand and eye tracking; \textbf{(3) Input devices}: motion controllers (Quest Touch, Index Knuckles), hand tracking, voice (microphone), bare-hand pinch (Vision Pro); \textbf{(4) Audio}: spatial audio with HRTF (binaural), often via on-strap speakers or earphones; \textbf{(5) Compute}: a tethered PC with GPU (PC-VR) or an on-board SoC (standalone); \textbf{(6) Software stack}: OS runtime (Quest OS, visionOS, SteamVR), \textbf{OpenXR runtime}, engine (Unity/Unreal/native), application logic; \textbf{(7) Optional}: haptics (vibration, gloves), treadmills, smell, biometric sensors. AR adds outward RGB+depth cameras and scene-understanding modules.

\switchcolumn

\textbf{Câu 25. Các thành phần chính của hệ thống VR/AR.}\\
Một hệ XR đầy đủ gồm: \textbf{(1) Display}: HMD (panel + lens); \textbf{(2) Tracking}: 6DoF head (IMU + camera), controller, hand, eye tracking; \textbf{(3) Input}: motion controller (Quest Touch, Index Knuckles), hand tracking, voice (mic), pinch (Vision Pro); \textbf{(4) Audio}: spatial audio HRTF (binaural) qua loa trên dây đeo hoặc tai nghe; \textbf{(5) Compute}: PC tethered có GPU (PC-VR) hoặc SoC on-board (standalone); \textbf{(6) Software stack}: OS runtime (Quest OS, visionOS, SteamVR), \textbf{OpenXR runtime}, engine (Unity/Unreal/native), logic ứng dụng; \textbf{(7) Tùy chọn}: haptic (rung, glove), treadmill, mùi, biometric. AR thêm camera RGB + depth hướng ngoài và module scene-understanding.

\switchcolumn*

\textbf{Q26. What is OpenXR and why is it important?}\\
\textbf{OpenXR} is an \textbf{open, royalty-free standard API} from the Khronos Group that lets a single application target many XR devices. Before OpenXR, developers wrote separate code for Oculus SDK, OpenVR, Windows Mixed Reality, Pico, etc.; now they call one API and the runtime (Meta, SteamVR, Pico, Snapdragon Spaces, visionOS via emulation) does the device-specific work. OpenXR exposes \textbf{spaces}, \textbf{actions}, \textbf{swapchains}, and \textbf{composition layers} in a vendor-neutral way. Importance: \textbf{(1) portability} (one codebase, many headsets); \textbf{(2) longevity} (apps survive HMD generations); \textbf{(3) ecosystem unity}; \textbf{(4) industry standard for enterprise procurement}. Unity (XR Plugin Management with OpenXR plugin) and Unreal natively target OpenXR; native C++/Vulkan engines link to it directly.

\switchcolumn

\textbf{Câu 26. OpenXR là gì? Vì sao quan trọng?}\\
\textbf{OpenXR} là \textbf{chuẩn API mở, miễn phí bản quyền} của Khronos Group, cho phép một app chạy trên nhiều thiết bị XR. Trước OpenXR phải viết riêng cho Oculus SDK, OpenVR, Windows MR, Pico... Giờ chỉ gọi một API, runtime (Meta, SteamVR, Pico, Snapdragon Spaces, visionOS qua emulation) lo phần thiết bị. OpenXR phơi ra \textbf{space}, \textbf{action}, \textbf{swapchain}, \textbf{composition layer} theo cách trung lập. Tầm quan trọng: \textbf{(1) portable} (một codebase, nhiều HMD); \textbf{(2) bền theo thời gian} (app sống qua nhiều thế hệ); \textbf{(3) thống nhất ecosystem}; \textbf{(4) chuẩn cho mua sắm doanh nghiệp}. Unity (XR Plugin Management + OpenXR plugin) và Unreal hỗ trợ OpenXR sẵn; engine C++/Vulkan link trực tiếp.

\switchcolumn*

\textbf{Q27. Compare ARKit, ARCore, and Vuforia.}\\
\textbf{ARKit} (Apple, iOS/visionOS): top-tier VI-SLAM, plane and mesh detection, body and face tracking, people occlusion, motion capture; iOS-only but rock-solid quality. \textbf{ARCore} (Google, Android + iOS): cross-platform mobile AR with motion tracking, plane finding, light estimation, \textbf{Cloud Anchors} for shared/persistent AR, Geospatial API (anchors at GPS+VPS coordinates worldwide). \textbf{Vuforia} (PTC): the veteran cross-platform SDK, strongest in \textbf{marker-based} (Image, Multi-Image, Cylinder, 3D Object Targets), Model Targets, Area Targets; works on iOS, Android, HoloLens, Magic Leap, Unity Editor. Choice: \textbf{Vuforia} for industrial / education / image-target apps; \textbf{ARKit} for premium iOS experiences; \textbf{ARCore} for Android-first, geo-located, or shared-anchor scenarios. Unity \textbf{AR Foundation} wraps ARKit + ARCore behind one API.

\switchcolumn

\textbf{Câu 27. So sánh ARKit, ARCore, Vuforia.}\\
\textbf{ARKit} (Apple, iOS/visionOS): VI-SLAM hàng đầu, plane + mesh, body + face tracking, people occlusion, mocap; chỉ iOS nhưng cực ổn. \textbf{ARCore} (Google, Android + iOS): mobile AR đa nền, motion tracking, plane, light estimation, \textbf{Cloud Anchors} cho AR chia sẻ/bền, Geospatial API (anchor theo GPS+VPS toàn cầu). \textbf{Vuforia} (PTC): SDK cross-platform lâu đời, mạnh nhất ở \textbf{marker-based} (Image, Multi-Image, Cylinder, 3D Object Target), Model Target, Area Target; chạy iOS, Android, HoloLens, Magic Leap, Unity Editor. Chọn: \textbf{Vuforia} cho công nghiệp/giáo dục/image-target; \textbf{ARKit} cho iOS cao cấp; \textbf{ARCore} cho Android-first, geo-located, shared anchor. Unity \textbf{AR Foundation} bọc ARKit + ARCore sau một API chung.

\switchcolumn*

\textbf{Q28. Describe the steps to build a VR app in Unity.}\\
\textbf{(1)} Install Unity LTS (e.g.\ 2022.3) with Android Build Support (for Quest); \textbf{(2)} Open Package Manager, install \textbf{XR Plugin Management}, then enable the \textbf{OpenXR} or \textbf{Oculus} provider; \textbf{(3)} Add the \textbf{XR Interaction Toolkit (XRI)} package for ray/grab/teleport interactors; \textbf{(4)} In the scene, replace the Main Camera with an \textbf{XR Origin (VR)} prefab (rig with tracked camera + two controllers); \textbf{(5)} Set up \textbf{tracked input actions} (Input System) for triggers, grips, thumbsticks; \textbf{(6)} Add \textbf{XR Grab Interactable} to objects, \textbf{Teleportation Areas} to floor; \textbf{(7)} Configure quality (URP, MSAA $\geq$~4x, 90~Hz); enable \textbf{single-pass instanced} rendering; \textbf{(8)} Build to Quest (Android) or PC (Windows + SteamVR/OpenXR); \textbf{(9)} Profile with the Oculus Performance HUD; iterate on draw calls, fill rate, and frame timing.

\switchcolumn

\textbf{Câu 28. Các bước build app VR trong Unity.}\\
\textbf{(1)} Cài Unity LTS (vd 2022.3) kèm Android Build Support (cho Quest); \textbf{(2)} Mở Package Manager, cài \textbf{XR Plugin Management}, bật provider \textbf{OpenXR} hoặc \textbf{Oculus}; \textbf{(3)} Cài \textbf{XR Interaction Toolkit (XRI)} cho ray/grab/teleport interactor; \textbf{(4)} Trong scene, thay Main Camera bằng prefab \textbf{XR Origin (VR)} (rig camera + 2 controller); \textbf{(5)} Setup \textbf{tracked input action} (Input System) cho trigger, grip, thumbstick; \textbf{(6)} Thêm \textbf{XR Grab Interactable} cho object, \textbf{Teleportation Area} cho sàn; \textbf{(7)} Cấu hình quality (URP, MSAA $\geq$~4x, 90~Hz), bật \textbf{single-pass instanced}; \textbf{(8)} Build cho Quest (Android) hoặc PC (Windows + SteamVR/OpenXR); \textbf{(9)} Profile bằng Oculus Performance HUD; tối ưu draw call, fill rate, frame time.

\switchcolumn*

\textbf{Q29. Describe the steps to build an AR app in Unity with AR Foundation.}\\
\textbf{(1)} Unity LTS + Build Support for iOS and Android; \textbf{(2)} Install \textbf{AR Foundation} + \textbf{ARKit XR Plugin} + \textbf{ARCore XR Plugin} via Package Manager; \textbf{(3)} Replace Main Camera with \textbf{AR Session Origin} (camera + tracked pose driver) and add an \textbf{AR Session} object; \textbf{(4)} Add managers for the features you need: \textbf{ARPlaneManager}, \textbf{ARRaycastManager}, \textbf{ARAnchorManager}, \textbf{ARFaceManager}, \textbf{ARMeshManager}, \textbf{AROcclusionManager}; \textbf{(5)} Write a script: tap $\to$ raycast against detected planes $\to$ instantiate a prefab anchored at the hit; \textbf{(6)} Configure player settings: ARM64, Vulkan/Metal, ARCore/ARKit required, camera permission strings; \textbf{(7)} Build and deploy to device (USB or TestFlight); \textbf{(8)} Iterate using \textbf{AR Foundation Editor Remote} or \textbf{XR Simulation} to test in editor without a phone.

\switchcolumn

\textbf{Câu 29. Các bước build app AR trong Unity với AR Foundation.}\\
\textbf{(1)} Unity LTS + Build Support cho iOS và Android; \textbf{(2)} Cài \textbf{AR Foundation} + \textbf{ARKit XR Plugin} + \textbf{ARCore XR Plugin} qua Package Manager; \textbf{(3)} Thay Main Camera bằng \textbf{AR Session Origin} (camera + tracked pose driver) và thêm object \textbf{AR Session}; \textbf{(4)} Thêm manager cho tính năng cần: \textbf{ARPlaneManager}, \textbf{ARRaycastManager}, \textbf{ARAnchorManager}, \textbf{ARFaceManager}, \textbf{ARMeshManager}, \textbf{AROcclusionManager}; \textbf{(5)} Viết script: chạm $\to$ raycast vào plane $\to$ instantiate prefab với anchor tại hit; \textbf{(6)} Player settings: ARM64, Vulkan/Metal, bắt buộc ARCore/ARKit, chuỗi camera permission; \textbf{(7)} Build và deploy lên thiết bị (USB hoặc TestFlight); \textbf{(8)} Lặp bằng \textbf{AR Foundation Editor Remote} hoặc \textbf{XR Simulation} để test trong editor.

\switchcolumn*

\textbf{Q30. Application: How is VR used in healthcare and surgical training?}\\
VR provides \textbf{safe, repeatable, high-fidelity} training environments for clinicians. (1) \textbf{Surgical simulators}: trainees rehearse laparoscopic, orthopedic, cataract, or endovascular procedures with haptic instruments (e.g.\ \textbf{Osso VR}, \textbf{FundamentalVR}, \textbf{Surgical Theater}, \textbf{ImmersiveTouch}); the system records metrics (path length, errors, time) for objective assessment. (2) \textbf{Anatomy education}: students explore segmented patient-specific 3D models reconstructed from CT/MRI (e.g.\ \textbf{3D Organon}, \textbf{Complete Anatomy}). (3) \textbf{Patient-specific surgical planning}: surgeon walks through the patient's CT in VR before operating. (4) \textbf{Pain distraction and anxiety reduction}: pediatric burns, dentistry, chemotherapy. \textbf{Tools}: Unity + XRI + Oculus/OpenXR, Vuforia for AR overlays, custom haptics (HaptX, Phantom Omni). The clinical evidence shows VR-trained residents achieve proficiency in significantly less OR time.

\switchcolumn

\textbf{Câu 30. Ứng dụng: VR trong y tế và đào tạo phẫu thuật.}\\
VR cung cấp môi trường huấn luyện \textbf{an toàn, lặp lại được, fidelity cao} cho bác sĩ. (1) \textbf{Simulator phẫu thuật}: học nội soi, chỉnh hình, đục thủy tinh thể, can thiệp mạch với dụng cụ haptic (\textbf{Osso VR}, \textbf{FundamentalVR}, \textbf{Surgical Theater}, \textbf{ImmersiveTouch}); hệ thống ghi metric (đường đi, lỗi, thời gian) để đánh giá khách quan. (2) \textbf{Giáo dục giải phẫu}: sinh viên khám phá mô hình 3D bệnh nhân từ CT/MRI (\textbf{3D Organon}, \textbf{Complete Anatomy}). (3) \textbf{Lập kế hoạch phẫu thuật cá nhân hóa}: bác sĩ ``đi qua'' CT bệnh nhân trong VR trước khi mổ. (4) \textbf{Giảm đau, giảm lo âu}: bỏng nhi, nha khoa, hóa trị. \textbf{Tool}: Unity + XRI + Oculus/OpenXR, Vuforia overlay, haptic riêng (HaptX, Phantom Omni). Bằng chứng lâm sàng: bác sĩ trẻ huấn luyện VR đạt thành thạo nhanh hơn nhiều.

\switchcolumn*

\textbf{Q31. Application: How is AR used in surgical navigation?}\\
\textbf{AR surgical navigation} overlays preoperative imaging (CT, MRI, vessels, tumor boundaries) directly onto the patient anatomy in the surgeon's view. Pipeline: \textbf{(1)} segment the CT/MRI to extract organs and target structures (3D Slicer, MITK); \textbf{(2)} register the preoperative model to the patient's body via \textbf{fiducial markers}, surface scanning, or feature-based registration; \textbf{(3)} track the patient and the instruments in real time (optical trackers, IR LEDs, or AR HMD inside-out); \textbf{(4)} render the segmented model at the registered pose on \textbf{HoloLens 2}, \textbf{Magic Leap 2}, or a tablet; \textbf{(5)} update with respiration / deformation. Real systems: \textbf{Brainlab Mixed Reality Viewer}, \textbf{Medivis SurgicalAR}, \textbf{Augmedics xvision} (FDA-approved for spine). Reported benefits: less radiation (fewer fluoroscopy shots), better screw placement accuracy, shorter OR time.

\switchcolumn

\textbf{Câu 31. Ứng dụng: AR trong dẫn đường phẫu thuật.}\\
\textbf{AR surgical navigation} chiếu hình ảnh tiền phẫu (CT, MRI, mạch máu, biên u) trực tiếp lên cơ thể bệnh nhân trong tầm nhìn bác sĩ. Pipeline: \textbf{(1)} segment CT/MRI để lấy cơ quan và cấu trúc đích (3D Slicer, MITK); \textbf{(2)} register mô hình tiền phẫu với cơ thể bằng \textbf{fiducial marker}, quét bề mặt, hoặc feature-based; \textbf{(3)} track bệnh nhân và dụng cụ theo thời gian thực (optical tracker, IR LED, hoặc AR HMD inside-out); \textbf{(4)} render mô hình ở pose đã register lên \textbf{HoloLens 2}, \textbf{Magic Leap 2}, hoặc tablet; \textbf{(5)} cập nhật theo hô hấp / biến dạng. Hệ thực: \textbf{Brainlab Mixed Reality Viewer}, \textbf{Medivis SurgicalAR}, \textbf{Augmedics xvision} (FDA cho cột sống). Lợi: giảm tia X, đặt vít chính xác hơn, rút ngắn thời gian mổ.

\switchcolumn*

\textbf{Q32. Application: VR/AR in education and training.}\\
VR/AR is transforming education by turning abstract concepts into \textbf{embodied, interactive 3D experiences}. \textbf{VR examples}: \textbf{Google Expeditions} / \textbf{Wonderscope} (virtual field trips); \textbf{Labster} (virtual labs --- run a chemistry experiment safely); \textbf{Engage}, \textbf{Mozilla Hubs}, \textbf{ENGAGE XR} for VR classrooms; \textbf{Tilt Brush} for art. \textbf{AR examples}: \textbf{JigSpace} for product visualization, \textbf{Merge Cube} for a holdable 3D model, \textbf{Quiver} for AR coloring books, \textbf{Civilisations AR} (BBC) for museum artifacts at home. \textbf{Why it works}: spatial presence improves recall; experiential learning matches Kolb's cycle; expensive/dangerous setups (chemistry, biology dissection, machinery) become safe and reusable. \textbf{Tools to build}: Unity + AR Foundation + Vuforia for image-target textbook AR; Unity + XRI for VR labs; \textbf{CoSpaces Edu} or \textbf{Spatial} for non-coders.

\switchcolumn

\textbf{Câu 32. Ứng dụng: VR/AR trong giáo dục và đào tạo.}\\
VR/AR biến khái niệm trừu tượng thành \textbf{trải nghiệm 3D nhập vai, tương tác}. \textbf{VR}: \textbf{Google Expeditions}/\textbf{Wonderscope} (đi tham quan ảo); \textbf{Labster} (lab ảo --- thí nghiệm hóa an toàn); \textbf{Engage}, \textbf{Mozilla Hubs}, \textbf{ENGAGE XR} làm lớp học VR; \textbf{Tilt Brush} cho mỹ thuật. \textbf{AR}: \textbf{JigSpace} cho product viz, \textbf{Merge Cube} cầm mô hình 3D, \textbf{Quiver} sách tô màu AR, \textbf{Civilisations AR} (BBC) đưa hiện vật bảo tàng về nhà. \textbf{Vì sao hiệu quả}: presence giúp nhớ lâu; học trải nghiệm theo Kolb; setup nguy hiểm/đắt (hóa, sinh, máy móc) trở nên an toàn và tái sử dụng. \textbf{Tool}: Unity + AR Foundation + Vuforia cho AR sách giáo khoa; Unity + XRI cho lab VR; \textbf{CoSpaces Edu} hoặc \textbf{Spatial} cho người không code.

\switchcolumn*

\textbf{Q33. Application: VR for cultural heritage, e-Museum, and e-Tourism.}\\
VR/AR enables \textbf{remote access} to fragile or distant heritage and increases \textbf{engagement} with cultural assets. Pipeline: scan the site/object using \textbf{photogrammetry} (Reality Capture, Meshroom), \textbf{LiDAR / terrestrial laser scanning}, or structured light; clean the mesh in Blender/MeshLab; texture it; import into Unity/Unreal; build a \textbf{VR walkthrough} (Quest) or \textbf{web AR / WebXR} viewer. Examples: \textbf{Mona Lisa: Beyond the Glass} (Louvre, HTC Vive), \textbf{Ancient Olympia in VR}, the digital reconstruction of \textbf{Notre-Dame de Paris}, \textbf{Smithsonian Voices VR}. For Vietnam: H{\`o}ang Th{\`a}nh Thăng Long, M\d{y} Sơn, Hu\^{e} Imperial City have all been targets of digital twin / VR projects (using NFD-based hole filling for damaged statues, Bezier surface modeling for restoration). Tools: Unity, Unreal, \textbf{8th Wall WebXR}, AR Foundation, Vuforia (image targets on museum cards).

\switchcolumn

\textbf{Câu 33. Ứng dụng: VR cho di sản văn hóa, e-Museum, e-Tourism.}\\
VR/AR cho phép \textbf{tiếp cận từ xa} các di sản dễ vỡ/xa xôi và tăng \textbf{tương tác} với tài sản văn hóa. Pipeline: quét địa điểm/hiện vật bằng \textbf{photogrammetry} (Reality Capture, Meshroom), \textbf{LiDAR/terrestrial laser scanning}, hoặc structured light; làm sạch mesh trong Blender/MeshLab; texture; nhập Unity/Unreal; làm \textbf{VR walkthrough} (Quest) hoặc \textbf{web AR/WebXR}. Ví dụ: \textbf{Mona Lisa: Beyond the Glass} (Louvre, Vive), \textbf{Olympia cổ đại VR}, dựng số \textbf{Notre-Dame Paris}, \textbf{Smithsonian Voices VR}. Việt Nam: Hoàng Thành Thăng Long, Mỹ Sơn, Cố đô Huế đều có dự án digital twin/VR (dùng NFD hole filling cho tượng vỡ, Bezier để phục dựng). Tool: Unity, Unreal, \textbf{8th Wall WebXR}, AR Foundation, Vuforia (image target trên thẻ bảo tàng).

\switchcolumn*

\textbf{Q34. Application: AR in industrial maintenance and Industry 4.0.}\\
\textbf{AR maintenance} overlays step-by-step procedures, IoT sensor data, and 3D animations onto the actual equipment, reducing errors and the need for senior experts on-site. Pipeline: \textbf{(1)} import the CAD model (CATIA/SolidWorks/Inventor) into Unity via \textbf{PiXYZ} or \textbf{CAD Exchanger}; \textbf{(2)} author procedures in \textbf{Vuforia Studio} or \textbf{Microsoft Dynamics 365 Guides}; \textbf{(3)} register the model to the real machine via \textbf{Vuforia Model Targets} or \textbf{Area Targets} (no marker needed --- the CAD itself is the recognition target); \textbf{(4)} run on \textbf{HoloLens 2}, \textbf{Magic Leap 2}, or rugged tablets; \textbf{(5)} stream live IoT data via \textbf{ThingWorx} / Azure Digital Twins; \textbf{(6)} \textbf{Remote Assist}: a senior expert sees the technician's view and annotates in real time. Documented benefits at Boeing, Airbus, Lockheed: 30--90\% time savings, fewer wiring errors.

\switchcolumn

\textbf{Câu 34. Ứng dụng: AR trong bảo trì công nghiệp và Industry 4.0.}\\
\textbf{AR maintenance} chiếu hướng dẫn từng bước, dữ liệu IoT, animation 3D lên máy thật $\Rightarrow$ giảm lỗi, ít cần chuyên gia tại chỗ. Pipeline: \textbf{(1)} đưa CAD (CATIA/SolidWorks/Inventor) vào Unity qua \textbf{PiXYZ} hoặc \textbf{CAD Exchanger}; \textbf{(2)} soạn quy trình trong \textbf{Vuforia Studio} hoặc \textbf{Dynamics 365 Guides}; \textbf{(3)} register mô hình với máy thật bằng \textbf{Vuforia Model Targets} hoặc \textbf{Area Targets} (không cần marker --- chính CAD là target); \textbf{(4)} chạy trên \textbf{HoloLens 2}, \textbf{Magic Leap 2}, hoặc tablet công nghiệp; \textbf{(5)} stream IoT qua \textbf{ThingWorx}/Azure Digital Twins; \textbf{(6)} \textbf{Remote Assist}: chuyên gia thấy view của thợ và chú thích thời gian thực. Boeing, Airbus, Lockheed: tiết kiệm 30--90\% thời gian, giảm lỗi.

\switchcolumn*

\textbf{Q35. Application: VR in psychotherapy --- exposure therapy and PTSD.}\\
\textbf{VR Exposure Therapy (VRET)} treats anxiety disorders, phobias, and PTSD by gradually exposing patients to controlled, virtual versions of feared stimuli (heights, flying, public speaking, combat). The therapist controls intensity, can pause, can repeat, and the patient never has to leave the clinic. Validated programs: \textbf{Bravemind} (USC ICT, for combat PTSD --- Iraq/Afghanistan/Vietnam scenarios), \textbf{Limbix VR}, \textbf{XRHealth} (FDA-cleared for several indications), \textbf{Psious / Amelia}. \textbf{Why it works}: emotional engagement is high (presence!); therapist can grade exposure precisely; cost is lower than real-world exposure (you cannot easily ``schedule a flight''). \textbf{Tools}: Unity + XRI + Oculus + biometric sensors (heart-rate variability, EDA, breathing) for biofeedback. Evidence base: dozens of RCTs show VRET non-inferior or superior to in-vivo exposure for several phobias.

\switchcolumn

\textbf{Câu 35. Ứng dụng: VR trong tâm lý trị liệu --- exposure therapy và PTSD.}\\
\textbf{VR Exposure Therapy (VRET)} điều trị rối loạn lo âu, ám ảnh sợ, PTSD bằng cách phơi nhiễm dần với phiên bản ảo của tác nhân sợ (cao, máy bay, nói trước đám đông, chiến đấu). Nhà trị liệu điều khiển cường độ, dừng, lặp; bệnh nhân không cần ra khỏi phòng khám. Chương trình đã kiểm chứng: \textbf{Bravemind} (USC ICT, PTSD chiến tranh --- Iraq/Afghanistan/Vietnam), \textbf{Limbix VR}, \textbf{XRHealth} (FDA-cleared), \textbf{Psious/Amelia}. \textbf{Vì sao hiệu quả}: cảm xúc cao (presence!); phân cấp phơi nhiễm chính xác; chi phí thấp hơn ngoài đời. \textbf{Tool}: Unity + XRI + Oculus + sensor sinh lý (HRV, EDA, hô hấp) cho biofeedback. Bằng chứng: hàng chục RCT cho thấy VRET không kém hoặc hơn phơi nhiễm thực với nhiều loại phobia.

\switchcolumn*

\textbf{Q36. Application: VR/AR in architecture, engineering, and construction (AEC).}\\
\textbf{AEC} firms use VR/AR throughout the design lifecycle. (1) \textbf{Design review in VR}: walk a 1:1 scale model of the building before construction (\textbf{Enscape}, \textbf{Twinmotion}, \textbf{IrisVR Prospect}, \textbf{Autodesk VRED}); detect ergonomic issues (low ceilings, cramped corridors) early. (2) \textbf{Client presentations}: photorealistic VR walkthroughs sell projects. (3) \textbf{On-site AR overlay}: surveyors view the BIM model registered to the construction site via HoloLens (\textbf{Trimble XR10}, \textbf{Trimble Connect}); they verify pipe placement, rebar layout. (4) \textbf{Clash detection} and coordination meetings in VR, multi-user. (5) \textbf{Safety training} (working at heights, hot work). \textbf{Pipeline}: Revit/ArchiCAD/Rhino $\to$ FBX/IFC $\to$ Unity Reflect or Unreal Datasmith $\to$ HMD. Real impact: fewer RFIs, fewer change orders, faster client buy-in.

\switchcolumn

\textbf{Câu 36. Ứng dụng: VR/AR trong kiến trúc và xây dựng (AEC).}\\
\textbf{AEC} dùng VR/AR xuyên suốt vòng đời thiết kế. (1) \textbf{Design review VR}: đi trong mô hình tỉ lệ 1:1 trước khi xây (\textbf{Enscape}, \textbf{Twinmotion}, \textbf{IrisVR Prospect}, \textbf{Autodesk VRED}); phát hiện trần thấp, hành lang chật từ sớm. (2) \textbf{Trình bày khách hàng}: walkthrough VR photorealistic. (3) \textbf{AR tại công trường}: kỹ sư xem BIM đè lên công trường qua HoloLens (\textbf{Trimble XR10}, \textbf{Trimble Connect}); kiểm tra ống, cốt thép. (4) \textbf{Clash detection} và họp phối hợp đa người trong VR. (5) \textbf{Đào tạo an toàn} (làm việc trên cao, hàn). \textbf{Pipeline}: Revit/ArchiCAD/Rhino $\to$ FBX/IFC $\to$ Unity Reflect hoặc Unreal Datasmith $\to$ HMD. Lợi ích thực: ít RFI, ít change order, khách duyệt nhanh.

\switchcolumn*

\textbf{Q37. Application: AR in retail and e-commerce.}\\
AR boosts conversion by letting customers \textbf{visualize products in their own space} before buying. (1) \textbf{Furniture try-out}: \textbf{IKEA Place}, \textbf{Wayfair View in Room}, \textbf{Amazon AR View} place a 3D sofa on the actual living-room floor with correct scale and shadows. (2) \textbf{AR try-on for fashion/cosmetics}: \textbf{Sephora Virtual Artist}, \textbf{L'Or{\'e}al ModiFace}, \textbf{Snapchat AR Lenses for Gucci/Dior}, eyeglass try-on (Warby Parker), shoe try-on (Nike). (3) \textbf{AR product packaging}: scan a soda can, see a campaign animation. \textbf{How it works}: ARKit/ARCore for plane detection and tracking; \textbf{ARFaceManager} for face mesh; \textbf{8th Wall} for browser AR (no install); \textbf{glTF} or \textbf{USDZ} as 3D format; \textbf{WebXR / Quick Look} delivery. \textbf{Documented impact}: AR-equipped product pages convert 2--3$\times$ better and reduce returns.

\switchcolumn

\textbf{Câu 37. Ứng dụng: AR trong bán lẻ và thương mại điện tử.}\\
AR tăng tỉ lệ chuyển đổi bằng cách cho khách \textbf{hình dung sản phẩm trong không gian riêng} trước khi mua. (1) \textbf{Thử nội thất}: \textbf{IKEA Place}, \textbf{Wayfair View in Room}, \textbf{Amazon AR View} đặt sofa 3D lên sàn phòng khách thật với tỉ lệ và bóng đúng. (2) \textbf{Thử thời trang/mỹ phẩm}: \textbf{Sephora Virtual Artist}, \textbf{L'Or{\'e}al ModiFace}, \textbf{Snapchat AR Lens cho Gucci/Dior}, thử kính (Warby Parker), thử giày (Nike). (3) \textbf{Bao bì AR}: quét lon nước ngọt $\to$ hiện animation chiến dịch. \textbf{Cách làm}: ARKit/ARCore cho plane + tracking; \textbf{ARFaceManager} cho face mesh; \textbf{8th Wall} cho AR trên browser (không cần cài); \textbf{glTF} hoặc \textbf{USDZ} làm format 3D; \textbf{WebXR/Quick Look} để phát. \textbf{Tác động}: trang có AR chuyển đổi cao gấp 2--3 lần và giảm tỉ lệ trả hàng.

\switchcolumn*

\textbf{Q38. Application: VR/AR in entertainment, social, and the metaverse.}\\
The largest commercial market for VR is \textbf{games and social experiences}. Hits: \textbf{Beat Saber}, \textbf{Half-Life: Alyx}, \textbf{Resident Evil 4 VR}, \textbf{Asgard's Wrath 2}, \textbf{Gorilla Tag}. Social platforms: \textbf{VRChat}, \textbf{Rec Room}, \textbf{Horizon Worlds}, \textbf{Spatial} support hundreds of users in shared virtual spaces with full-body avatars, lip-sync, and \textbf{eye contact} via eye tracking. AR side: \textbf{Pokemon GO}, \textbf{Snapchat Lenses}, \textbf{TikTok AR effects} reach billions. Production tools: Unity + Photon Fusion / \textbf{Mirror} / \textbf{Normcore} for networking; \textbf{Ready Player Me} avatars; \textbf{Lens Studio} (Snap), \textbf{Spark AR} (Meta), \textbf{Effect House} (TikTok) for AR filters. Concert/live: \textbf{Travis Scott in Fortnite}, \textbf{ABBA Voyage} mixed-reality stage. The ``metaverse'' vision unifies these: persistent, multi-user, cross-device 3D internet.

\switchcolumn

\textbf{Câu 38. Ứng dụng: VR/AR trong giải trí, mạng xã hội và metaverse.}\\
Thị trường thương mại lớn nhất của VR là \textbf{game và mạng xã hội}. Hit: \textbf{Beat Saber}, \textbf{Half-Life: Alyx}, \textbf{Resident Evil 4 VR}, \textbf{Asgard's Wrath 2}, \textbf{Gorilla Tag}. Nền tảng xã hội: \textbf{VRChat}, \textbf{Rec Room}, \textbf{Horizon Worlds}, \textbf{Spatial} hỗ trợ hàng trăm người trong không gian ảo, full-body avatar, lip-sync, \textbf{giao tiếp mắt} qua eye tracking. AR: \textbf{Pokemon GO}, \textbf{Snapchat Lens}, \textbf{TikTok AR effect} tiếp cận hàng tỉ người. Tool: Unity + Photon Fusion/\textbf{Mirror}/\textbf{Normcore} cho networking; \textbf{Ready Player Me} cho avatar; \textbf{Lens Studio} (Snap), \textbf{Spark AR} (Meta), \textbf{Effect House} (TikTok) cho AR filter. Sự kiện trực tiếp: \textbf{Travis Scott trong Fortnite}, \textbf{ABBA Voyage}. Tầm nhìn ``metaverse'': internet 3D bền, đa người dùng, đa thiết bị.

\switchcolumn*

\textbf{Q39. What sensors are needed for AR tracking? (and what each one provides)}\\
A modern AR device combines complementary sensors: \textbf{(1) Wide-angle RGB camera(s)} for visual features, image targets, hand and face tracking; \textbf{(2) IMU} (3-axis gyroscope + 3-axis accelerometer + sometimes magnetometer) at 800--1000~Hz for fast rotational tracking and motion prediction; \textbf{(3) Depth sensor} (\textbf{ToF} on Pixel/iPad LiDAR, or \textbf{structured light} on iPhone Pro front for FaceID, or \textbf{stereo} pairs) for plane detection, occlusion, mesh reconstruction; \textbf{(4) GPS + barometer} for outdoor coarse localization; \textbf{(5) Magnetometer} for heading; \textbf{(6) Microphones} for voice + spatial audio. \textbf{Sensor fusion} (Extended Kalman Filter / factor graphs) combines these into a single, smooth 6DoF pose. The IMU provides high frequency and short-term accuracy; the camera corrects long-term drift; depth gives metric scale and occlusion.

\switchcolumn

\textbf{Câu 39. AR tracking cần sensor nào? (mỗi loại làm gì)}\\
Thiết bị AR hiện đại phối hợp các sensor bù trừ: \textbf{(1) Camera RGB góc rộng} cho feature, image target, hand + face; \textbf{(2) IMU} (gyro 3 trục + accel 3 trục + đôi khi magnetometer) 800--1000~Hz cho xoay nhanh và motion prediction; \textbf{(3) Depth sensor} (\textbf{ToF} trên Pixel/iPad LiDAR, hoặc \textbf{structured light} trên iPhone Pro mặt trước cho FaceID, hoặc \textbf{stereo}) cho plane, occlusion, dựng mesh; \textbf{(4) GPS + barometer} cho định vị ngoài trời; \textbf{(5) Magnetometer} cho hướng; \textbf{(6) Mic} cho voice + spatial audio. \textbf{Sensor fusion} (EKF / factor graph) gộp lại thành 6DoF pose mượt. IMU cho tần số cao + chính xác ngắn hạn; camera sửa drift dài hạn; depth cho scale chuẩn và occlusion.

\switchcolumn*

\textbf{Q40. What are the main challenges in VR? How are they mitigated?}\\
\textbf{(1) Cybersickness}: caused by latency or artificial locomotion --- mitigate with $\geq$~90~Hz, M2P $<$~20~ms, teleport, snap turn, vignette. \textbf{(2) Hardware ergonomics}: weight on the face, heat, sweat, cable --- mitigated by lighter pancake optics, balanced strap (Bigscreen Beyond, Quest Elite), wireless. \textbf{(3) Eye fatigue / VAC}: mismatch between vergence and fixed lens accommodation --- research on \textbf{varifocal} (Half Dome, Vision Pro hints), \textbf{light-field} displays. \textbf{(4) Content cost}: 3D assets, animation, optimization --- partly addressed by asset stores, photogrammetry, AI-generated content. \textbf{(5) Onboarding} for non-gamers: tutorials, hand tracking instead of controllers. \textbf{(6) Compute budget} on standalone --- foveated rendering, fixed-foveation, level-of-detail, baked lighting. \textbf{(7) Social acceptance / privacy} --- guardian boundaries, cameras off in private spaces.

\switchcolumn

\textbf{Câu 40. Thách thức chính của VR và cách giảm thiểu.}\\
\textbf{(1) Cybersickness}: do latency hoặc locomotion nhân tạo --- giải bằng $\geq$~90~Hz, M2P $<$~20~ms, teleport, snap turn, vignette. \textbf{(2) Ergonomic}: nặng mặt, nóng, mồ hôi, dây --- pancake optics nhẹ, dây cân bằng (Bigscreen Beyond, Quest Elite), wireless. \textbf{(3) Mỏi mắt/VAC}: lệch vergence vs accommodation --- nghiên cứu \textbf{varifocal} (Half Dome, gợi ý Vision Pro), \textbf{light-field}. \textbf{(4) Chi phí nội dung}: asset 3D, animation, tối ưu --- asset store, photogrammetry, AI-generated. \textbf{(5) Onboarding} cho người không phải gamer: tutorial, hand tracking thay controller. \textbf{(6) Compute} trên standalone --- foveated rendering, fixed foveation, LOD, baked lighting. \textbf{(7) Xã hội/quyền riêng tư} --- guardian, camera tắt trong không gian riêng.

\switchcolumn*

\textbf{Q41. What are the main challenges in AR (drift, occlusion, lighting, registration)?}\\
\textbf{(1) Tracking drift}: visual--inertial SLAM accumulates error; mitigated by \textbf{loop closure}, \textbf{cloud anchors}, persistent maps, fiducial markers as anchors. \textbf{(2) Occlusion}: virtual objects must hide behind real ones --- solved with \textbf{depth maps} (ARKit People Occlusion, ARCore Depth API), \textbf{scene mesh}, or instance segmentation. \textbf{(3) Lighting consistency}: virtuals must match real light --- \textbf{Light Estimation} (ambient intensity, color temperature, spherical harmonics, primary directional source) provided by ARKit/ARCore; HDR environment probes. \textbf{(4) Registration accuracy}: alignment between virtual model and real object --- improved with marker-based or model-target tracking and recalibration. \textbf{(5) Display limits}: small FOV (HoloLens 52$^\circ$), low brightness in sunlight, semi-transparent virtuals on OST. \textbf{(6) Power and thermals} on phones / glasses --- duty cycling, low-power tracking modes.

\switchcolumn

\textbf{Câu 41. Thách thức chính của AR (drift, occlusion, ánh sáng, registration).}\\
\textbf{(1) Drift của tracking}: VI-SLAM cộng dồn sai số --- giải bằng \textbf{loop closure}, \textbf{cloud anchor}, bản đồ bền, dùng marker làm anchor. \textbf{(2) Occlusion}: vật ảo phải lẩn sau vật thật --- dùng \textbf{depth map} (ARKit People Occlusion, ARCore Depth API), \textbf{scene mesh}, hoặc instance segmentation. \textbf{(3) Ánh sáng nhất quán}: vật ảo phải hợp ánh sáng thật --- \textbf{Light Estimation} (cường độ, color temperature, spherical harmonics, hướng chính) của ARKit/ARCore; HDR env probe. \textbf{(4) Độ chính xác registration}: căn chỉnh model ảo với vật thật --- dùng marker hoặc model target và recalibrate. \textbf{(5) Hạn chế display}: FOV nhỏ (HoloLens 52$^\circ$), độ sáng thấp ngoài nắng, vật ảo trong mờ trên OST. \textbf{(6) Pin và nhiệt} trên phone/kính --- duty cycle, low-power tracking.

\switchcolumn*

\textbf{Q42. Why does AR need plane detection, and how does it work?}\\
Without \textbf{plane detection}, an AR app cannot place virtual objects on real surfaces with correct scale and orientation; objects would float arbitrarily. Pipeline (ARCore/ARKit): \textbf{(1)} The SLAM module continuously produces a \textbf{point cloud} of feature points; \textbf{(2)} Points are clustered by spatial proximity; \textbf{(3)} \textbf{RANSAC} fits planes through clusters with low residual error; \textbf{(4)} Each plane is classified as \textbf{horizontal} (floor, table) or \textbf{vertical} (wall) using gravity from the IMU; \textbf{(5)} Plane boundaries are estimated using a convex hull or alpha shape; \textbf{(6)} Planes grow over frames as more points are seen; nearby planes can merge. The app uses \textbf{ARRaycastManager} to cast a ray from a screen tap to the planes and instantiate a prefab at the hit. On LiDAR-equipped iPads, the depth map is used directly for instant plane detection and meshing.

\switchcolumn

\textbf{Câu 42. Vì sao AR cần plane detection và hoạt động ra sao?}\\
Không có \textbf{plane detection}, app AR không đặt được vật ảo lên bề mặt thật đúng tỉ lệ + hướng; vật sẽ trôi lung tung. Pipeline (ARCore/ARKit): \textbf{(1)} SLAM sinh \textbf{point cloud} feature liên tục; \textbf{(2)} Cluster point theo khoảng cách; \textbf{(3)} \textbf{RANSAC} fit plane qua cluster với sai số nhỏ; \textbf{(4)} Plane được phân loại \textbf{ngang} (sàn, bàn) hay \textbf{dọc} (tường) dựa vào gravity từ IMU; \textbf{(5)} Biên plane ước lượng bằng convex hull hoặc alpha shape; \textbf{(6)} Plane lớn dần qua các frame; plane gần nhau có thể merge. App dùng \textbf{ARRaycastManager} bắn tia từ điểm chạm tới plane và instantiate prefab tại hit. Trên iPad có LiDAR, depth map dùng trực tiếp để dò plane và mesh tức thì.

\switchcolumn*

\textbf{Q43. Application: VR/AR in cultural and tourism --- e-Tourism Vietnam case.}\\
\textbf{e-Tourism} brings physical destinations to remote audiences and enhances on-site visits. (1) \textbf{360$^\circ$ VR videos / VR walkthroughs} of Ha Long Bay, Phong Nha caves, Tr{\`a}ng An, Cu Chi tunnels --- captured with \textbf{Insta360 Pro 2}, served via YouTube VR or a Unity native app on Quest. (2) \textbf{AR guides at heritage sites}: scan a QR code or use VPS (ARCore Geospatial API) to overlay reconstructions of original temple roofs, painted murals, or Cham towers onto today's ruins. (3) \textbf{Photogrammetric digital twins}: scan statues with a mobile phone (Polycam, RealityScan), reconstruct online, and let visitors examine details impossible to see from a distance. (4) \textbf{Hole filling \& restoration}: damaged statues at My Son Sanctuary can be ``virtually restored'' using normal-field-driven hole filling (the topic of NFD-HoleFill) and Bezier surface modeling. \textbf{Tools}: Unity, Reality Composer, AR Foundation, 8th Wall WebXR for in-browser AR (no install needed for tourists).

\switchcolumn

\textbf{Câu 43. Ứng dụng: VR/AR trong văn hóa và du lịch --- e-Tourism Việt Nam.}\\
\textbf{e-Tourism} đưa địa điểm thực tới khán giả từ xa và làm phong phú trải nghiệm tại chỗ. (1) \textbf{Video VR 360$^\circ$/VR walkthrough} Hạ Long, Phong Nha, Tràng An, địa đạo Củ Chi --- quay bằng \textbf{Insta360 Pro 2}, phát qua YouTube VR hoặc app Unity trên Quest. (2) \textbf{AR guide tại di tích}: quét QR hoặc dùng VPS (ARCore Geospatial API) để overlay phục dựng mái đền, tranh tường, tháp Chăm lên di tích hiện tại. (3) \textbf{Digital twin photogrammetry}: quét tượng bằng phone (Polycam, RealityScan), dựng online, cho khách xem chi tiết không thấy từ xa. (4) \textbf{Hole filling và phục dựng}: tượng vỡ ở Mỹ Sơn có thể ``phục chế ảo'' bằng normal-field-driven hole filling (đề tài NFD-HoleFill) và Bezier surface. \textbf{Tool}: Unity, Reality Composer, AR Foundation, 8th Wall WebXR cho AR trong browser (du khách không cần cài app).

\switchcolumn*

\textbf{Q44. What are the key performance metrics for VR/AR (latency, FOV, refresh rate, presence)?}\\
\textbf{(1) Motion-to-photon latency} ($<$~20~ms target); \textbf{(2) Frame rate / refresh rate} (90~Hz minimum, 120/144~Hz preferred); \textbf{(3) Resolution} (per eye, e.g.\ Quest 3 2064$\times$2208) and \textbf{pixels per degree (PPD)} ($>$~25 reduces SDE; Vision Pro $\sim$34); \textbf{(4) FOV} (110$^\circ$ for VR, 50$^\circ$ for HoloLens AR); \textbf{(5) Tracking accuracy and jitter} (sub-millimeter for Lighthouse, $\sim$mm for inside-out); \textbf{(6) Tracking volume} (room-scale 3$\times$3 m typical); \textbf{(7) Battery / weight / thermals} on standalone; \textbf{(8) Subjective: presence questionnaires} (SUS, Witmer--Singer); \textbf{(9) Cybersickness scores} (Simulator Sickness Questionnaire SSQ); \textbf{(10) Task performance} (time, errors) for training and clinical apps. \textbf{(11) For AR}: registration accuracy in mm, drift over time, light estimation error, FPS under thermal load.

\switchcolumn

\textbf{Câu 44. Các metric hiệu năng quan trọng của VR/AR.}\\
\textbf{(1) Motion-to-photon latency} (mục tiêu $<$~20~ms); \textbf{(2) FPS/refresh} (tối thiểu 90~Hz, tốt 120/144~Hz); \textbf{(3) Độ phân giải} mỗi mắt (Quest 3 2064$\times$2208) và \textbf{pixel per degree (PPD)} ($>$~25 giảm SDE; Vision Pro $\sim$34); \textbf{(4) FOV} (110$^\circ$ VR, 50$^\circ$ HoloLens); \textbf{(5) Độ chính xác và jitter của tracking} (sub-mm Lighthouse, $\sim$mm inside-out); \textbf{(6) Vùng tracking} (room-scale 3$\times$3 m); \textbf{(7) Pin/nặng/nhiệt} trên standalone; \textbf{(8) Chủ quan: questionnaire presence} (SUS, Witmer--Singer); \textbf{(9) Cybersickness} (SSQ); \textbf{(10) Hiệu suất task} (thời gian, lỗi) cho training và y tế. \textbf{(11) AR}: độ chính xác registration mm, drift theo thời gian, sai số light estimation, FPS dưới tải nhiệt.

\switchcolumn*

\textbf{Q45. What are future trends in VR/AR?}\\
\textbf{(1) Light, all-day AR glasses} (Meta Orion prototype, Apple AR glasses, Snap Spectacles 5) replacing the smartphone over the next decade. \textbf{(2) Pancake / metalens optics} for thin, high-clarity HMDs. \textbf{(3) Eye- and face-tracking everywhere} for foveated rendering, social presence, and biometric input. \textbf{(4) Higher resolution} ($>$~4K per eye, retina-class); micro-OLED mainstream. \textbf{(5) AI-generated 3D content}: text-to-3D (Luma AI Genie, Meshy, Stable Dreamer), one-shot avatars, neural radiance fields for capture. \textbf{(6) Generative agents and LLM-driven NPCs} in social VR. \textbf{(7) WebXR and \textbf{8th Wall} growing browser-based AR}. \textbf{(8) Standardization} (OpenXR, glTF, USD/USDZ). \textbf{(9) Healthcare and enterprise verticals} maturing (FDA-cleared therapeutics). \textbf{(10) Brain-computer interfaces} and neural wristbands for input (Meta CTRL-Labs). \textbf{(11) Spatial computing} (Apple's term) merges VR/AR into general-purpose desktop replacement.

\switchcolumn

\textbf{Câu 45. Xu hướng tương lai của VR/AR.}\\
\textbf{(1) Kính AR nhẹ, đeo cả ngày} (Meta Orion, Apple AR glasses, Snap Spectacles 5) thay smartphone trong thập kỷ tới. \textbf{(2) Pancake/metalens optics} cho HMD mỏng, sắc nét. \textbf{(3) Eye + face tracking khắp nơi} cho foveated rendering, social presence, biometric input. \textbf{(4) Resolution cao hơn} ($>$~4K mỗi mắt, retina-class); micro-OLED phổ thông. \textbf{(5) AI sinh nội dung 3D}: text-to-3D (Luma AI Genie, Meshy, Stable Dreamer), avatar one-shot, NeRF cho capture. \textbf{(6) Generative agent + LLM-driven NPC} trong social VR. \textbf{(7) WebXR và \textbf{8th Wall}} đẩy AR trên browser. \textbf{(8) Chuẩn hóa} (OpenXR, glTF, USD/USDZ). \textbf{(9) Y tế và doanh nghiệp} trưởng thành (therapeutic được FDA duyệt). \textbf{(10) BCI / dây đeo cổ tay thần kinh} làm input (Meta CTRL-Labs). \textbf{(11) Spatial computing} (Apple) hợp nhất VR/AR thành máy tính cá nhân thế hệ mới.

\switchcolumn*

\textbf{Q46. Exam-style: ``Application of VR in [field]'' --- a template answer you can adapt.}\\
\textbf{Field}: [healthcare / education / culture / industry / retail / psychology / AEC / military].\\
Structure your essay in 5 parts:\\
\textbf{(1) Problem statement}: what real-world pain (training cost, danger, accessibility, scale, knowledge transfer) does VR solve here? \textbf{(2) How VR helps}: which specific affordance (presence, safe failure, visualization, repetition, remote presence) addresses the problem? \textbf{(3) How it works} (3--5 steps): capture data $\to$ build 3D scene $\to$ track user $\to$ render stereo $\to$ interact + assess. \textbf{(4) Tools \& techniques}: HMD (Quest 3 / HoloLens 2 / Vision Pro), engine (Unity + XR Interaction Toolkit), SDKs (OpenXR, ARKit/ARCore, Vuforia for marker), supporting tech (photogrammetry, MRI segmentation, BIM, IoT, biometrics). \textbf{(5) Concrete example} (a real product or paper): name a system, the institution, and a measured outcome. End with one sentence on remaining limitations. This 5-part structure handles any ``application'' question on the exam.

\switchcolumn

\textbf{Câu 46. Mẫu trả lời ``Ứng dụng VR trong [lĩnh vực]'' có thể tái sử dụng.}\\
\textbf{Lĩnh vực}: [y tế / giáo dục / văn hóa / công nghiệp / bán lẻ / tâm lý / AEC / quân sự].\\
Cấu trúc 5 phần:\\
\textbf{(1) Vấn đề}: nỗi đau thực tế nào (chi phí huấn luyện, nguy hiểm, khó tiếp cận, quy mô, truyền tri thức) mà VR giải được? \textbf{(2) VR giúp thế nào}: affordance cụ thể (presence, sai mà không nguy hiểm, trực quan, lặp, hiện diện từ xa) trị vấn đề ra sao? \textbf{(3) Cách hoạt động} (3--5 bước): thu dữ liệu $\to$ dựng scene 3D $\to$ track người dùng $\to$ render stereo $\to$ tương tác + đánh giá. \textbf{(4) Tool \& kỹ thuật}: HMD (Quest 3/HoloLens 2/Vision Pro), engine (Unity + XRI), SDK (OpenXR, ARKit/ARCore, Vuforia cho marker), công nghệ phụ (photogrammetry, segmentation MRI, BIM, IoT, biometric). \textbf{(5) Ví dụ cụ thể}: tên hệ thống, đơn vị, kết quả đo được. Kết bằng một câu về giới hạn còn lại. Cấu trúc 5 phần này dùng được cho mọi câu ``ứng dụng'' trong đề thi.

\end{paracol}
